
\subsection{Derechos patrimoniales}
\label{sec:derechos-patrimoniales}

El presente trabajo de grado es desarrollado por Andrés Felipe Pinzón Márquez, Daniel Felipe Ramírez Vargas, Guillermo Andrés Aponte Cárdenas y Miguel Ángel Sarmiento Cuervo, estudiantes de Ingeniería de Sistemas de la Pontificia Universidad Javeriana, bajo la dirección académica de Jaime Andrés Pavlich Mariscal.

El proyecto comprende la creación de una aplicación de escritorio, código fuente, documentación técnica, componentes de consulta y visualización, configuraciones de simulación y una biblioteca de escenarios precalculados.

\textbf{Titularidad de los desarrollos originales.}

Para efectos de esta propuesta, y considerando su carácter de trabajo académico desarrollado por iniciativa de los estudiantes, se establece como régimen de titularidad propuesto que los derechos patrimoniales sobre las obras originales creadas por el equipo correspondan a sus cuatro integrantes, en calidad de autores o coautores según sus aportes creativos.

Esta declaración se fundamenta en las Directrices de Propiedad Intelectual de la Pontificia Universidad Javeriana, aprobadas en 2018, que reconocen la titularidad de los estudiantes sobre determinadas obras académicas desarrolladas por iniciativa propia y bajo su cuenta y riesgo. También contemplan las obras realizadas bajo dirección de un profesor cuando esta no implique su participación directa en la creación \cite{vfpJaverianaDirectrices2018}.

La dirección académica de Jaime Andrés Pavlich Mariscal se reconoce como parte fundamental del acompañamiento del proyecto. Dicha función no implica automáticamente cotitularidad patrimonial, sin perjuicio del reconocimiento de autoría que corresponda si existieran aportes creativos directos a los resultados protegibles.

El software se encuentra amparado por las normas colombianas de derecho de autor. El Decreto 1360 de 1989 reconoce el soporte lógico como una creación protegida y comprende dentro de este concepto el programa de computador, su descripción y el material auxiliar \cite{vfpDecreto1360}.

La titularidad aquí propuesta presupone que el proyecto no constituye una obra realizada por encargo institucional ni se encuentra sujeto a acuerdos previos de cesión o cotitularidad. Si se estableciera una situación distinta, como una investigación contratada, cofinanciada o desarrollada bajo condiciones de uso significativo de recursos universitarios que afecten la titularidad, se aplicarán las disposiciones y acuerdos correspondientes.

\textbf{Autorización académica a la Universidad.}

Los estudiantes conservarán los derechos patrimoniales sobre sus desarrollos originales bajo las condiciones anteriormente señaladas y prevén otorgar a la Pontificia Universidad Javeriana una autorización de uso del documento final del trabajo de grado para su conservación, consulta y difusión académica mediante la Biblioteca y el repositorio institucional.

De acuerdo con el numeral 21 de las Directrices de Propiedad Intelectual, esta autorización puede otorgarse sin efectuar una cesión de los derechos patrimoniales \cite{vfpJaverianaDirectrices2018}. Su formalización se realizará mediante los instrumentos institucionales correspondientes.

La autorización de consulta del documento no supone, por sí sola, la entrega de derechos de modificación, distribución comercial o explotación del código fuente de la aplicación. Cualquier autorización adicional sobre el software deberá establecer expresamente su objeto, alcance y condiciones.

\textbf{Herramientas y fuentes de información de terceros.}

El proyecto utilizará Fire Dynamics Simulator (FDS) y Smokeview para la simulación y visualización de resultados. Para el procesamiento geográfico se contempla el uso de QGIS y el complemento qgis2fds, que permite preparar entradas territoriales para FDS y procesar determinadas salidas georreferenciadas \cite{fdsnist,vfpQGIS2FDS}.

Entre las fuentes de información identificadas se encuentran modelos digitales de elevación como GEDTM30 y Copernicus DEM GLO-30; información meteorológica de ERA5 y ERA5-Land; y datos complementarios provenientes del IDEAM, los portales de datos abiertos de Bogotá y otras fuentes territoriales pertinentes.

La incorporación efectiva de estas fuentes dependerá de su disponibilidad, resolución, pertinencia y compatibilidad con el dominio seleccionado. Se respetarán las licencias, requisitos de atribución y condiciones de reutilización aplicables a cada herramienta y conjunto de datos \cite{vfpCopernicusDEM,vfpCopernicusERA5}.

El empleo de software de terceros, información geográfica y recursos computacionales universitarios no implica atribuir al equipo la propiedad intelectual sobre dichos elementos.

\textbf{Alcance de uso de la aplicación.}

La finalidad principal del producto será apoyar al comandante o coordinador de una emergencia en la consulta e interpretación de escenarios de posible propagación de incendios forestales.

De manera complementaria, la aplicación podrá utilizarse para explorar escenarios en actividades académicas, demostraciones y ejercicios de capacitación. Este uso secundario no constituye un servicio de formación contratado por un cuerpo de bomberos ni representa una transferencia de titularidad a las instituciones que eventualmente participen en las pruebas.

La aplicación no reemplazará el criterio profesional del personal de emergencia ni se presentará como una herramienta operativa validada sin las comprobaciones correspondientes.

En consecuencia, la titularidad patrimonial propuesta se mantiene en cabeza de los estudiantes creadores, sin perjuicio de los derechos de terceros, las autorizaciones académicas expresamente otorgadas y las condiciones institucionales o contractuales que resulten aplicables.
